
@ID: OW22D1P1I
@BIDANG: Analisis Real

Diketahui fungsi $f : \mathbb{R} \mapsto \mathbb{R}$ memenuhi $|f(x) - f(y)| \le |x^3 - x^2y|^2$, untuk setiap $x, y \in \mathbb{R}$. Jika $f(1) = 1$, maka $f(2022) = \ldots$

@ID: OW22D1P2I
@BIDANG: Analisis Real

Jika untuk setiap $n \in \mathbb{N}$, $a_n = \sum_{k=0}^{n} \frac{1}{n+k}$, maka nilai $\lim_{n \to \infty} a_n = \ldots$

@ID: OW22D1P3I
@BIDANG: Struktur Aljabar

Misalkan $\sigma = \begin{pmatrix} 1 & 2 & 3 & 4 & 5 \\ 3 & 4 & 1 & 5 & 2 \end{pmatrix}$ dan $\tau = \begin{pmatrix} 1 & 2 & 3 & 4 & 5 \\ 5 & 2 & 1 & 3 & 4 \end{pmatrix}$. Baris kedua dari 

$$
\alpha = \begin{pmatrix} 1 & 2 & 3 & 4 & 5 \\ \ldots & \ldots & \ldots & \ldots & \ldots \end{pmatrix}
$$

agar $\alpha = (\sigma \circ \tau)^{-1}$ adalah $\ldots$

@ID: OW22D1P4I
@BIDANG: Struktur Aljabar

Semua $y \in \frac{\mathbb{Z}_8[x]}{\langle x^2 - 1 \rangle}$ sehingga $(7 + 5x)y = 0$ adalah $\ldots$

@ID: OW22D1P5I
@BIDANG: Kombinatorika

Misalkan $n$ adalah bilangan asli yang terletak diantara 1 dan 2022 yang memuat digit 0. Banyaknya bilangan $n$ yang demikian adalah $\ldots$

---

@ID: OW22D1P1U
@BIDANG: Analisis Real

Diberikan fungsi $f : [1, \infty) \to \mathbb{R}$ memenuhi $f(1) = 1$ dan 

$$
f'(x) = \frac{x}{x^4 + [f(x)]^4}
$$

untuk setiap $x \in [1, \infty)$. Buktikan bahwa $\lim_{x \to \infty} f(x)$ ada dengan nilainya tidak melebihi $1 + \frac{\pi}{8}$.

@ID: OW22D1P2U
@BIDANG: Analisis Real

Diberikan barisan bilangan real $X = (x_n)$, dengan $x_1 = \sqrt{2}$ dan $x_n = (\sqrt{2})^{x_{n-1}}$ untuk $n > 1$. Buktikan bahwa $X$ konvergen, kemudian tentukan nilai limitnya.

@ID: OW22D1P3U
@BIDANG: Struktur Aljabar

Diberikan grup $G$ dan $\phi : G \to G$ adalah suatu homomorfisma grup. Buktikan atau sangkal pernyataan berikut: $\phi$ adalah suatu isomorfisma jika dan hanya jika untuk setiap $g \in G$, orde $g$ sama dengan orde $\phi(g)$.

@ID: OW22D1P4U
@BIDANG: Struktur Aljabar

Misalkan $R = \frac{\mathbb{Z}_p[x]}{\langle x^2 - x \rangle}$, dimana $p$ adalah bilangan prima. Tunjukkan bahwa 

$$
S = \{ (1-x)a \mid a \in R \}
$$

isomorfik dengan $\mathbb{Z}_p$.

@ID: OW22D1P5U
@BIDANG: Kombinatorika

Misalkan $n$ adalah suatu bilangan 12 digit yang disusun dari 3,4,5,6, atau 7. Jika setiap digit pada $n$ muncul paling sedikit 2 kali dan paling banyak 4 kali, tentukan banyaknya susunan bilangan $n$ yang dapat dibentuk.
@ID: OW22D2P1I
@BIDANG: Aljabar Linear

Jika $A$ matriks berukuran $2\times2$ dengan $\det(A)>2$ yang memenuhi persamaan

$$
A^2=2A+\begin{pmatrix}3&10\\0&8\end{pmatrix},
$$

maka $\det(A)=\ldots$

@ID: OW22D2P2I
@BIDANG: Aljabar Linear

Misalkan $P_1$ adalah ruang polinom dengan koefisien real berderajat paling tinggi $1$ yang dilengkapi dengan hasil kali dalam

$$
\langle f,g\rangle=\int_0^1 f(x)g(x)\,dx
$$

untuk setiap $f,g\in P_1$.

Diberikan pemetaan linear $T:P_1\to\mathbb R^3$ yang memenuhi

$$
T(1)=\begin{pmatrix}1\\2\\4\end{pmatrix}
\quad\text{dan}\quad
T(2x-1)=\begin{pmatrix}3\\6\\12\end{pmatrix}.
$$

Jika polinom $h(x)=ax+b$ dengan $b\neq0$ tegak lurus terhadap setiap vektor di

$$
\{\,f\in P_1:T(f)=0\,\},
$$

maka $\dfrac{a}{b}=\ldots$

@ID: OW22D2P3I
@BIDANG: Analisis Kompleks

Banyak bilangan kompleks tak real $z$ yang memenuhi

$$
z^{20}=22
$$

adalah $\ldots$

@ID: OW22D2P4I
@BIDANG: Analisis Kompleks

Jika

$$
z=\sin\left(\frac{\pi}{4}+i\right),
$$

maka

$$
\sqrt2\,(\operatorname{Re}(z)+\operatorname{Im}(z))
=\ldots
$$

(Tuliskan jawaban dalam bentuk paling sederhana.)

@ID: OW22D2P5I
@BIDANG: Kombinatorika

Sisi-sisi dari petak persegi satuan sebuah papan catur $4\times4$ diwarnai merah. Banyaknya sisi satuan (panjangnya $1$ satuan) minimal yang harus diwarnai sehingga setiap petak persegi satuan memiliki paling sedikit tiga sisi berwarna merah adalah $\ldots$

@ID: OW22D2P1U
@BIDANG: Aljabar Linear

Tentukan semua pasangan bilangan real $a$ dan $b$ sehingga sistem persamaan linear

$$
\begin{pmatrix}
a&1&0\\
0&a&1\\
-b&0&a
\end{pmatrix}
\begin{pmatrix}
x\\
y\\
z
\end{pmatrix}
=
\begin{pmatrix}
1\\
2\\
-3
\end{pmatrix}
$$

konsisten (memiliki paling sedikit satu solusi). Jelaskan jawaban Anda.

@ID: OW22D2P2U
@BIDANG: Aljabar Linear

Misalkan $A=(a_{ij})$ dan $B=(b_{ij})$ matriks berukuran $n\times n$ dengan entri bilangan real. Jika

$$
\operatorname{tr}(AB)
=
\frac12
\sum_{i=1}^n
\sum_{j=1}^n
(a_{ij}^2+b_{ij}^2),
$$

buktikan bahwa $A^T=B$.

@ID: OW22D2P3U
@BIDANG: Analisis Kompleks

Diberikan bilangan-bilangan kompleks $a$ dan $b$ dengan $|a|=|b|=1$. Tunjukkan bahwa

$$
z=(a+b)\left(\frac1a+\frac1b\right)
$$

merupakan bilangan real dan memenuhi

$$
0\le z\le4.
$$

@ID: OW22D2P4U
@BIDANG: Analisis Kompleks

Diberikan fungsi kompleks

$$
f(z)=
\begin{cases}
\dfrac{z^3}{z^2}, & \text{jika } z\neq0,\\
0, & \text{jika } z=0.
\end{cases}
$$

Apakah fungsi $f$ terdiferensialkan kompleks di titik $0$? Jelaskan!

@ID: OW22D2P5U
@BIDANG: Kombinatorika

Diberikan $30$ titik pada koordinat Kartesius, yaitu titik $D_i(i,1)$ untuk $i\in\{1,2,\ldots,15\}$ dan titik $D_j(j-15,4)$ untuk $j\in\{16,17,\ldots,30\}$.

Tentukan banyaknya segitiga sama kaki yang dapat dibentuk dengan ketiga titik sudutnya diambil dari $\{D_1,D_2,\ldots,D_{30}\}$.